\scanpage{24}
On notera cependant que la donnée d'une
forme $\varphi$ \inserted{sur $M$} est \uline{plus} \uline{précis} que la donnée de
l'involution sur $A=\mathrm{End}(M)$ ; ainsi si $\mathrm{End}(M)$
est commutatif, ttes les $\varphi$ donnent la
$\deleted{\unsure{m}}=$ involution ! De $=$, si $\lambda\in Q^*$, $\lambda\varphi$ donne
la $=$ involution que $\varphi$

Considérons la forme \uline{symétrique}
\[
\mathrm{Tr}\,v^\varphi u=\mathrm{Tr}\,u^\varphi v
\]
sur \inserted{le motif} \deleted{$A=$} $\underline{\mathrm{End}}(M)\simeq$ \deleted{\unsure{induite}, par la forme} \inserted{$M^\vee\otimes M$}
cette forme n'est autre que $\overset{\vee}{\varphi}\otimes\varphi$ où
$\overset{\vee}{\varphi}$ est la forme \uline{duale} de $\varphi$ sur $M^\vee$, \inserted{déduite de $\varphi$ par transport de structure par $\widetilde\varphi$} \deleted{\unsure{ou}}
\deleted{\unsure{auquel} on \unsure{a} identifié $\underline{\mathrm{End}}(M)$ à
$\unsure{\Phi}(M)\simeq M^\vee\otimes M^\vee$ et \uline{supposons} qu'elle soit
dé}
Elle induit une forme sur
$\mathrm{End}(M)$ à valeurs ds $Q$, \unsure{notre point|notre forme|motif point|notée point}.

\uline{Supposons} que cette forme soit \deleted{\uline{définie}}
\uline{positive} (elle est néc. non dég, donc \uline{définie}
positive). \deleted{Soit $P$ le cône des} Alors
$A=\mathrm{End}(M)$ devient une ``algèbre de Riemann''.
% archive sheet number: 23
